\section{Introduction}
\label{sec:intro}

Why can ``how to think generally'' be formalized at all? General thinking---
moving from a concrete solution to a cost model to a unified form, checking that
the original solution is a special case, and landing every qualitative claim on
a testable number---looks like a craft. We argue that it has a stable internal
structure that admits a \emph{descriptive} formalization: a small set of
principles from which the observed practices follow as consequences. Our
contribution is descriptive; it is not a claim about the mechanisms of human or
artificial cognition.

We make three contributions. \emph{First}, a core formalization: understanding
as a lossless round-trip fixed point of a compression--decomposition pair
($\Dmap\circ\Cmap=\mathrm{id}$), with losslessness defined as uniqueness plus
reproducibility. \emph{Second}, a consistency test: we derive seven
thinking-style hypotheses from the framework and obtain 7/7 internal consistency
on the source corpus (an internal consistency check, not external validation).
\emph{Third}, an actionable protocol: five injectable prompt templates
(T1--T5), each specifying a trigger condition and an action, positioned as an
agent-prompt injection protocol.

\paragraph{Scope and honesty.}
This paper is a descriptive formalization plus an actionable protocol. We do
\emph{not} claim to have discovered the mechanism of human or AI thinking, nor
do we offer cognitive-science conclusions. The consistency test is
self-referential: the framework was constructed from and tested on the same
corpus. Template effectiveness is pending independent experiments; we claim
injectability and operability only. All quantitative statements (7/7, five
templates, four facets) are traceable to the source corpus (see the appendix and
the accompanying number-trace file).
